\item ~\defpoints{25} [Boosting]

Suppose that we are interested in learning a classifier, such that at any turn of a game we can pose a question, like ``should I attack this ant hill now?", and get an answer.That is, we want to build a classifier which we can feed some features on the current game state, and get the output ``attack" or ``don't attack". There are many possible ways to define what the action ``attack" means, but for now let's define it as sending all friendly ants that can see the ant hill under consideration towards it.

Let's recall the AdaBoost algorithm described in class. Its input is a dataset $\{(x_{i},y_{i})\}_{i=1}^{n}$, with $x_i$ being the $i$-th sample, and $y_{i}\in \{-1,1\}$ denoting the $i$-th label, $i=1,2,...,n$. The features might be composed of a count of the number of friendly ants that can see the ant hill under consideration, and a count of the number of enemy ants these friendly ants can see. For example, if there were 10 friendly ants that could see a particular ant hill, and 5 enemy ants that the friendly ants could see, we would have:
$$x_1 = \begin{bmatrix} 10 \\ 5 \end{bmatrix}$$

The label of the example $x_{1}$ is $y_{1} = 1$, once the friendly ants were successful in razing the enemy ant hill, and $y_{1} = -1$ otherwise. We could generate such examples by running a greedy bot (or any other opponent bot) against a bot that we periodically try to attack an enemy ant hill. Each time this bot tries the attack, we record (say, after $20$ turns or some other significant amount of time) whether the attack was successful or not.

\begin{itemize}
\item[(a)] Let $\epsilon_t$ denote the error of a weak classifier $h_t$:
$$\epsilon_{t} = \sum_{i=1}^{n} D_{t}(i) \mathbb{I}(y_{i} \neq h_{t}(x_{i}))$$
In the simple “attack” / “don't attack” scenario, suppose that we have implemented the following six weak classifiers:
\begin{align*}
h^{(1)}(x_{i}) = 2 * \mathbb{I}(x_{i1} \geq 2) - 1, \hspace{1cm}  & h^{(4)}(x_{i}) = 2 * \mathbb{I}(x_{i2} \leq 2) - 1,  \\
h^{(2)}(x_{i}) = 2 * \mathbb{I}(x_{i1} \geq 6) - 1, \hspace{1cm}  & h^{(5)}(x_{i}) = 2 * \mathbb{I}(x_{i2} \leq 6) - 1,  \\
h^{(3)}(x_{i}) = 2 * \mathbb{I}(x_{i1} \geq 10) - 1, \hspace{1cm} & h^{(6)}(x_{i}) = 2 * \mathbb{I}(x_{i2} \leq 10) - 1.
\end{align*}
Given ten training data points ($n = 10$) as shown in Table \ref{table1}:
\begin{table}[htbp]
    \centering
    \begin{tabular}{|c|cc|c|}
    \hline
    $i$ & $x_{i1}$ & $x_{i2}$ & $y_{i}$ \\ \hline
    1 & 1.5 & 0.5 & 1 \\
    2 & 2.5 & 1.5 & 1 \\
    3 & 3.5 & 3.5 & 1 \\
    4 & 6.5 & 5.5 & 1 \\
    5 & 7.5 & 10.5 & 1 \\
    6 & 1.5 & 2.5 & -1 \\
    7 & 3.5 & 1.5 & -1 \\
    8 & 5.5 & 5.5 & -1 \\
    9 & 7.5 & 8.5 & -1 \\
    10 & 1.5 & 10.5 & -1 \\
    \hline
    \end{tabular}
    \caption{The training data in (a).}
    \label{table1}
\end{table}

please show that what is the minimum value of $\epsilon_{1}$ and which of $h^{(1)},\cdots ,h^{(6)}$ achieve this value? Note that there may be multiple classifiers that all have the same $\epsilon_{1}$. You should list all classifiers that achieve the minimum $\epsilon_{1}$ value. ~\defpoints{5}

\item[(b)] For all the questions in the remainder of this section, let $h_{1}$ denote $h^{(1)}$ chosen in the first round of boosting. (That is, $h^{(1)}$ was the classifier that achieved the minimum $\epsilon_{1}$.)
\begin{enumerate}
    \item[(1)] What is the value of $\alpha_{1}$ (the weight of this first classifier $h_{1}$)? ~\defpoints{2}

    \item[(2)] What should $Z_{t}$ be in order to make sure the distribution $D_{t+1}$ is normalized correctly? That is, derive the formula of $Z_{t}$ in terms of $\epsilon_{t}$ that will ensure $\sum\limits_{i=1}^{n} D_{t+1}(i) = 1$. Please also derive the formula of $\alpha_{t}$ in terms of $\epsilon_{t}$. ~\defpoints{5}

    \item[(3)] Which points will increase in significance in the second round of boosting? That is, for which points will we have $D_{1}(i) < D_{2}(i)$? What are the values of $D_{2}$ for these points? ~\defpoints{5}

    \item[(4)] In the second round of boosting, the weights on the points will be different, and thus the error $\epsilon_2$ will also be different. Which of $h^{(1)}, \cdots, h^{(6)}$ will minimize $\epsilon_2$? (Which classifier will be selected as the second weak classifier $h_2$?) What is its value of $\epsilon_2$? ~\defpoints{5}

    \item[(5)] What will the average error of the final classifier $H$ be, if we stop after these two rounds of boosting? That is, if $H(x) = \text{sign}(\alpha_{1}h_{1}(x) + \alpha_{2}h_{2}(x))$, what will the  training error $\epsilon = \dfrac{1}{n} \sum\limits_{i=1}^{n} \mathbb{I} (y_{i} \neq H(x_{i}))$ be? Is this more, less, or the same as the error we would get, if we just used one of the weak classifiers instead of this final classifier $H$ ~\defpoints{3}

\end{enumerate}
\end{itemize}

\solution
\subsection*{(a) Minimum $\epsilon_1$ and corresponding classifiers}
The initial weights $D_1(i)$ for each data point $i=1, \ldots, 10$ are uniform: $D_1(i) = \frac{1}{n} = \frac{1}{10} = 0.1$.
The error of a weak classifier $h_t$ is given by $\epsilon_t = \sum_{i=1}^{n} D_t(i) \mathbb{I}(y_i \neq h_t(x_i))$.
Since $D_1(i) = 0.1$, $\epsilon_1 = 0.1 \times (\text{number of misclassified points})$.
The six weak classifiers are:
\begin{align*}
h^{(1)}(x_{i}) &= 2 \cdot \mathbb{I}(x_{i1} \geq 2) - 1 \\
h^{(2)}(x_{i}) &= 2 \cdot \mathbb{I}(x_{i1} \geq 6) - 1 \\
h^{(3)}(x_{i}) &= 2 \cdot \mathbb{I}(x_{i1} \geq 10) - 1 \\
h^{(4)}(x_{i}) &= 2 \cdot \mathbb{I}(x_{i2} \leq 2) - 1 \\
h^{(5)}(x_{i}) &= 2 \cdot \mathbb{I}(x_{i2} \leq 6) - 1 \\
h^{(6)}(x_{i}) &= 2 \cdot \mathbb{I}(x_{i2} \leq 10) - 1
\end{align*}
Let's evaluate the number of misclassifications for each classifier using the data from Table \ref{table1} in the problem statement.

\begin{itemize}
    \item For $h^{(1)}(x_i) = 2 \cdot \mathbb{I}(x_{i1} \geq 2) - 1$:
    \begin{itemize}
        \item $x_1=(1.5, 0.5), y_1=1$: $x_{11}=1.5 < 2 \Rightarrow h^{(1)}(x_1)=-1$. Misclassified.
        \item $x_2=(2.5, 1.5), y_2=1$: $x_{21}=2.5 \geq 2 \Rightarrow h^{(1)}(x_2)=1$. Correct.
        \item $x_3=(3.5, 3.5), y_3=1$: $x_{31}=3.5 \geq 2 \Rightarrow h^{(1)}(x_3)=1$. Correct.
        \item $x_4=(6.5, 5.5), y_4=1$: $x_{41}=6.5 \geq 2 \Rightarrow h^{(1)}(x_4)=1$. Correct.
        \item $x_5=(7.5, 10.5), y_5=1$: $x_{51}=7.5 \geq 2 \Rightarrow h^{(1)}(x_5)=1$. Correct.
        \item $x_6=(1.5, 2.5), y_6=-1$: $x_{61}=1.5 < 2 \Rightarrow h^{(1)}(x_6)=-1$. Correct.
        \item $x_7=(3.5, 1.5), y_7=-1$: $x_{71}=3.5 \geq 2 \Rightarrow h^{(1)}(x_7)=1$. Misclassified.
        \item $x_8=(5.5, 5.5), y_8=-1$: $x_{81}=5.5 \geq 2 \Rightarrow h^{(1)}(x_8)=1$. Misclassified.
        \item $x_9=(7.5, 8.5), y_9=-1$: $x_{91}=7.5 \geq 2 \Rightarrow h^{(1)}(x_9)=1$. Misclassified.
        \item $x_{10}=(1.5, 10.5), y_{10}=-1$: $x_{10,1}=1.5 < 2 \Rightarrow h^{(1)}(x_{10})=-1$. Correct.
    \end{itemize}
    Misclassified points for $h^{(1)}$: $\{1, 7, 8, 9\}$. Number of misclassifications = 4. $\epsilon_1^{(1)} = 0.1 \times 4 = 0.4$.

    \item For $h^{(2)}(x_i) = 2 \cdot \mathbb{I}(x_{i1} \geq 6) - 1$:
    Misclassified points: $\{1, 2, 3, 9\}$. Number of misclassifications = 4. $\epsilon_1^{(2)} = 0.1 \times 4 = 0.4$.

    \item For $h^{(3)}(x_i) = 2 \cdot \mathbb{I}(x_{i1} \geq 10) - 1$:
    Misclassified points: $\{1, 2, 3, 4, 5\}$. Number of misclassifications = 5. $\epsilon_1^{(3)} = 0.1 \times 5 = 0.5$.

    \item For $h^{(4)}(x_i) = 2 \cdot \mathbb{I}(x_{i2} \leq 2) - 1$:
    Misclassified points: $\{3, 4, 5, 7\}$. Number of misclassifications = 4. $\epsilon_1^{(4)} = 0.1 \times 4 = 0.4$.

    \item For $h^{(5)}(x_i) = 2 \cdot \mathbb{I}(x_{i2} \leq 6) - 1$:
    Misclassified points: $\{5, 6, 7, 8\}$. Number of misclassifications = 4. $\epsilon_1^{(5)} = 0.1 \times 4 = 0.4$.

    \item For $h^{(6)}(x_i) = 2 \cdot \mathbb{I}(x_{i2} \leq 10) - 1$:
    Misclassified points: $\{5, 6, 7, 8, 9\}$. Number of misclassifications = 5. $\epsilon_1^{(6)} = 0.1 \times 5 = 0.5$.
\end{itemize}
The minimum value of $\epsilon_1$ is $0.4$.
The classifiers that achieve this minimum value are $h^{(1)}$, $h^{(2)}$, $h^{(4)}$, and $h^{(5)}$.

\subsection*{(b) Assuming $h_1 = h^{(1)}$}
We are given that $h_1 = h^{(1)}$, so $\epsilon_1 = 0.4$.

\subsubsection*{(1) Value of $\alpha_1$}
The weight $\alpha_t$ of a classifier $h_t$ is given by:
$$ \alpha_t = \frac{1}{2} \ln\left(\frac{1-\epsilon_t}{\epsilon_t}\right) $$
For $t=1$, $\epsilon_1 = 0.4$:
$$ \alpha_1 = \frac{1}{2} \ln\left(\frac{1-0.4}{0.4}\right) = \frac{1}{2} \ln\left(\frac{0.6}{0.4}\right) = \frac{1}{2} \ln(1.5) $$
Numerically, $\ln(1.5) \approx 0.405465$, so $\alpha_1 \approx 0.20273$.

\subsubsection*{(2) Derivation of $Z_t$ and $\alpha_t$}
The distribution $D_{t+1}$ is updated as follows:
$$ D_{t+1}(i) = \frac{D_t(i) \exp(-\alpha_t y_i h_t(x_i))}{Z_t} $$
To ensure $\sum_{i=1}^{n} D_{t+1}(i) = 1$, $Z_t$ must be the sum of the numerators:
$$ Z_t = \sum_{i=1}^{n} D_t(i) \exp(-\alpha_t y_i h_t(x_i)) $$
Let $C_t$ be the set of indices of correctly classified samples by $h_t$ ($y_i h_t(x_i)=1$), and $M_t$ be the set of indices of misclassified samples ($y_i h_t(x_i)=-1$).
Then $Z_t$ can be written as:
$$ Z_t = \sum_{i \in C_t} D_t(i) e^{-\alpha_t} + \sum_{i \in M_t} D_t(i) e^{\alpha_t} $$
Since $\epsilon_t = \sum_{i \in M_t} D_t(i)$ and $\sum_{i=1}^n D_t(i) = 1$, it follows that $\sum_{i \in C_t} D_t(i) = 1 - \epsilon_t$.
So, $Z_t$ in terms of $\epsilon_t$ and $\alpha_t$ is:
$$ Z_t = (1-\epsilon_t)e^{-\alpha_t} + \epsilon_t e^{\alpha_t} $$
The value of $\alpha_t$ is chosen to minimize $Z_t$. We differentiate $Z_t$ with respect to $\alpha_t$ and set it to zero:
$$ \frac{dZ_t}{d\alpha_t} = -(1-\epsilon_t)e^{-\alpha_t} + \epsilon_t e^{\alpha_t} = 0 $$
$$ \epsilon_t e^{\alpha_t} = (1-\epsilon_t)e^{-\alpha_t} $$
$$ e^{2\alpha_t} = \frac{1-\epsilon_t}{\epsilon_t} $$
$$ \alpha_t = \frac{1}{2} \ln\left(\frac{1-\epsilon_t}{\epsilon_t}\right) $$
This is the formula for $\alpha_t$. Substituting this $\alpha_t$ back into the expression for $Z_t$:
$e^{\alpha_t} = \sqrt{\frac{1-\epsilon_t}{\epsilon_t}}$ and $e^{-\alpha_t} = \sqrt{\frac{\epsilon_t}{1-\epsilon_t}}$.
$$ Z_t = (1-\epsilon_t)\sqrt{\frac{\epsilon_t}{1-\epsilon_t}} + \epsilon_t\sqrt{\frac{1-\epsilon_t}{\epsilon_t}} = \sqrt{(1-\epsilon_t)\epsilon_t} + \sqrt{\epsilon_t(1-\epsilon_t)} = 2\sqrt{\epsilon_t(1-\epsilon_t)} $$
This is the formula for $Z_t$ in terms of $\epsilon_t$.

\subsubsection*{(3) Points with increased significance and their $D_2$ values}
A point $i$ will increase in significance if $D_2(i) > D_1(i)$.
$D_2(i) = D_1(i) \frac{\exp(-\alpha_1 y_i h_1(x_i))}{Z_1}$.
This occurs if $\exp(-\alpha_1 y_i h_1(x_i)) / Z_1 > 1$.
Since $\alpha_1 = \frac{1}{2}\ln(1.5) > 0$, $\exp(\alpha_1) > 1$ and $\exp(-\alpha_1) < 1$.
$Z_1 = 2\sqrt{\epsilon_1(1-\epsilon_1)} = 2\sqrt{0.4(0.6)} = 2\sqrt{0.24} \approx 2 \times 0.4899 \approx 0.9798$.
Since $Z_1 < 1$:
\begin{itemize}
    \item If $y_i h_1(x_i) = -1$ (point $i$ is misclassified by $h_1$), then $\exp(-\alpha_1 y_i h_1(x_i)) = e^{\alpha_1} = \sqrt{1.5} \approx 1.2247$. Then $e^{\alpha_1}/Z_1 \approx 1.2247/0.9798 \approx 1.25 > 1$. So $D_2(i) > D_1(i)$.
    \item If $y_i h_1(x_i) = 1$ (point $i$ is correctly classified by $h_1$), then $\exp(-\alpha_1 y_i h_1(x_i)) = e^{-\alpha_1} = 1/\sqrt{1.5} \approx 0.8165$. Then $e^{-\alpha_1}/Z_1 \approx 0.8165/0.9798 \approx 0.833 < 1$. So $D_2(i) < D_1(i)$.
\end{itemize}
The points that increase in significance are those misclassified by $h_1 = h^{(1)}$. From part (a), these are points $\{1, 7, 8, 9\}$.
For these misclassified points, $D_1(i) = 0.1$:
$$ D_2(i) = \frac{D_1(i) e^{\alpha_1}}{Z_1} = \frac{0.1 \sqrt{(1-\epsilon_1)/\epsilon_1}}{2\sqrt{\epsilon_1(1-\epsilon_1)}} = \frac{0.1}{2\epsilon_1} = \frac{0.1}{2 \times 0.4} = \frac{0.1}{0.8} = 0.125 $$
So, for points $i \in \{1, 7, 8, 9\}$, $D_2(i) = 0.125$.
For correctly classified points $i \in \{2, 3, 4, 5, 6, 10\}$:
$$ D_2(i) = \frac{D_1(i) e^{-\alpha_1}}{Z_1} = \frac{0.1 \sqrt{\epsilon_1/(1-\epsilon_1)}}{2\sqrt{\epsilon_1(1-\epsilon_1)}} = \frac{0.1}{2(1-\epsilon_1)} = \frac{0.1}{2 \times 0.6} = \frac{0.1}{1.2} = \frac{1}{12} \approx 0.08333 $$
Check sum: $4 \times 0.125 + 6 \times \frac{1}{12} = 0.5 + 0.5 = 1$. The distribution $D_2$ is normalized.

\subsubsection*{(4) Second weak classifier $h_2$ and its error $\epsilon_2$}
We calculate $\epsilon_2^{(j)} = \sum_{i=1}^{10} D_2(i) \mathbb{I}(y_i \neq h^{(j)}(x_i))$ for each $h^{(j)}$.
Recall $D_2(i)=0.125$ for $i \in \{1,7,8,9\}$ and $D_2(i)=1/12$ for $i \in \{2,3,4,5,6,10\}$.

\begin{itemize}
    \item $h^{(1)}$: Misclassified set $M_1^{(1)}=\{1,7,8,9\}$.
    $\epsilon_2^{(1)} = D_2(1)+D_2(7)+D_2(8)+D_2(9) = 4 \times 0.125 = 0.5$.
    \item $h^{(2)}$: Misclassified set $M_1^{(2)}=\{1,2,3,9\}$.
    $\epsilon_2^{(2)} = D_2(1)+D_2(2)+D_2(3)+D_2(9) = 0.125 + (1/12) + (1/12) + 0.125 = 0.25 + 2/12 = 1/4 + 1/6 = 3/12 + 2/12 = 5/12 \approx 0.4167$.
    My previous calculation was: $0.25 + 1/60 = 16/60 = 4/15$. Checking again. $D_2(2)=1/12$. $0.25 + 2/12 = 0.25 + 1/6 = 1/4 + 1/6 = (3+2)/12 = 5/12$. Ah, the previous $1/60$ was wrong. $0.25 + 0.01666...$ was wrong. $1/6 \approx 0.1666...$. So $0.25 + 0.1666... = 0.41666...$.
    $\epsilon_2^{(2)} = 0.125 (\text{for } 1) + (1/12) (\text{for } 2) + (1/12) (\text{for } 3) + 0.125 (\text{for } 9) = 0.25 + 2/12 = 1/4 + 1/6 = 3/12 + 2/12 = 5/12$.

    Let's recalculate all $\epsilon_2^{(j)}$ values systematically:
    $S_w = \{1,7,8,9\}$ (weight $0.125 = 1/8$)
    $S_c = \{2,3,4,5,6,10\}$ (weight $1/12$)

    $\epsilon_2^{(1)}$: $\{1,7,8,9\}$. All in $S_w$. $\epsilon_2^{(1)} = 4 \times (1/8) = 1/2 = 0.5$.
    $\epsilon_2^{(2)}$: $\{1,2,3,9\}$. $1,9 \in S_w$; $2,3 \in S_c$. $\epsilon_2^{(2)} = 2 \times (1/8) + 2 \times (1/12) = 1/4 + 1/6 = 3/12 + 2/12 = 5/12 \approx 0.4167$.
    $\epsilon_2^{(3)}$: $\{1,2,3,4,5\}$. $1 \in S_w$; $2,3,4,5 \in S_c$. $\epsilon_2^{(3)} = 1 \times (1/8) + 4 \times (1/12) = 1/8 + 1/3 = 3/24 + 8/24 = 11/24 \approx 0.4583$.
    $\epsilon_2^{(4)}$: $\{3,4,5,7\}$. $7 \in S_w$; $3,4,5 \in S_c$. $\epsilon_2^{(4)} = 1 \times (1/8) + 3 \times (1/12) = 1/8 + 1/4 = 1/8 + 2/8 = 3/8 = 0.375$.
    $\epsilon_2^{(5)}$: $\{5,6,7,8\}$. $7,8 \in S_w$; $5,6 \in S_c$. $\epsilon_2^{(5)} = 2 \times (1/8) + 2 \times (1/12) = 1/4 + 1/6 = 5/12 \approx 0.4167$.
    $\epsilon_2^{(6)}$: $\{5,6,7,8,9\}$. $7,8,9 \in S_w$; $5,6 \in S_c$. $\epsilon_2^{(6)} = 3 \times (1/8) + 2 \times (1/12) = 3/8 + 1/6 = 9/24 + 4/24 = 13/24 \approx 0.5417$.
\end{itemize}
Comparing the $\epsilon_2^{(j)}$ values:
$\epsilon_2^{(1)} = 0.5$
$\epsilon_2^{(2)} = 5/12 \approx 0.4167$
$\epsilon_2^{(3)} = 11/24 \approx 0.4583$
$\epsilon_2^{(4)} = 3/8 = 0.375$
$\epsilon_2^{(5)} = 5/12 \approx 0.4167$
$\epsilon_2^{(6)} = 13/24 \approx 0.5417$

The minimum value is $\epsilon_2 = 3/8 = 0.375$. This is achieved by $h^{(4)}$.
So, the second weak classifier selected is $h_2 = h^{(4)}$, and its weighted error is $\epsilon_2 = 3/8$.

\subsubsection*{(5) Training error of the final classifier $H(x)$}
The final classifier after two rounds is $H(x) = \operatorname{sign}(\alpha_1 h_1(x) + \alpha_2 h_2(x))$.
We have $h_1 = h^{(1)}$ with $\epsilon_1 = 0.4$, so $\alpha_1 = \frac{1}{2}\ln(1.5) \approx 0.20273$.
And $h_2 = h^{(4)}$ with $\epsilon_2 = 3/8 = 0.375$.
$\alpha_2 = \frac{1}{2} \ln\left(\frac{1-\epsilon_2}{\epsilon_2}\right) = \frac{1}{2} \ln\left(\frac{1-3/8}{3/8}\right) = \frac{1}{2} \ln\left(\frac{5/8}{3/8}\right) = \frac{1}{2} \ln(5/3)$.
Numerically, $\ln(5/3) \approx \ln(1.6667) \approx 0.5108$, so $\alpha_2 \approx 0.2554$.

We evaluate $H(x_i) = \operatorname{sign}(\alpha_1 h^{(1)}(x_i) + \alpha_2 h^{(4)}(x_i))$ for each $x_i$.
It's equivalent to evaluate $\operatorname{sign}(\ln(1.5) h^{(1)}(x_i) + \ln(5/3) h^{(4)}(x_i))$.
$\ln(1.5) \approx 0.405465$. $\ln(5/3) \approx 0.510826$.

\begin{table}[htbp]
\centering
\begin{tabular}{|c|c|c|c|c|c|p{3.5cm}|c|c|}
\hline
$i$ & $x_{i1}$ & $x_{i2}$ & $y_i$ & $h^{(1)}(x_i)$ & $h^{(4)}(x_i)$ & $\ln(1.5)h^{(1)} + \ln(5/3)h^{(4)}$ & $H(x_i)$ & $y_i \neq H(x_i)$? \\ \hline
1 & 1.5 & 0.5 & 1 & -1 & 1 & $-0.405465(1) + 0.510826(1) = 0.105361$ & 1 & No \\
2 & 2.5 & 1.5 & 1 & 1 & 1 & $0.405465(1) + 0.510826(1) = 0.916291$ & 1 & No \\
3 & 3.5 & 3.5 & 1 & 1 & -1 & $0.405465(1) + 0.510826(-1) = -0.105361$ & -1 & Yes \\
4 & 6.5 & 5.5 & 1 & 1 & -1 & $0.405465(1) + 0.510826(-1) = -0.105361$ & -1 & Yes \\
5 & 7.5 & 10.5 & 1 & 1 & -1 & $0.405465(1) + 0.510826(-1) = -0.105361$ & -1 & Yes \\
6 & 1.5 & 2.5 & -1 & -1 & -1 & $0.405465(-1) + 0.510826(-1) = -0.916291$ & -1 & No \\
7 & 3.5 & 1.5 & -1 & 1 & 1 & $0.405465(1) + 0.510826(1) = 0.916291$ & 1 & Yes \\
8 & 5.5 & 5.5 & -1 & 1 & -1 & $0.405465(1) + 0.510826(-1) = -0.105361$ & -1 & No \\
9 & 7.5 & 8.5 & -1 & 1 & -1 & $0.405465(1) + 0.510826(-1) = -0.105361$ & -1 & No \\
10 & 1.5 & 10.5 & -1 & -1 & -1 & $0.405465(-1) + 0.510826(-1) = -0.916291$ & -1 & No \\
\hline
\end{tabular}
\caption{Evaluation of $H(x_i)$}
\label{tab:H_eval}
\end{table}
The misclassified points for $H(x)$ are $\{3, 4, 5, 7\}$. There are 4 misclassifications.
The training error is $\epsilon = \frac{1}{n} \sum_{i=1}^{n} \mathbb{I}(y_i \neq H(x_i)) = \frac{4}{10} = 0.4$.

Comparing this to the errors of the weak classifiers (with uniform weights, as in part (a)):
Error of $h_1 = h^{(1)}$ was $0.4$.
Error of $h_2 = h^{(4)}$ was $0.4$.
The error of the final classifier $H(x)$ ($0.4$) is the same as the error of these individual weak classifiers.

\newpage